\documentclass[10pt,a4paper]{amsart}
\usepackage[margin=19mm]{geometry}
\usepackage{amsmath,amssymb,mathtools}
\usepackage[scr=rsfs,cal=euler]{mathalfa}
\usepackage{xfrac}
\usepackage[unicode,hidelinks,bookmarks=false]{hyperref}
\newcommand{\ie}{i.e.\ }
\newcommand{\cf}{cf.\ }
\newcommand{\resp}{respectively\ }
\newcommand{\Subsection}{\subsection}
\providecommand{\nobreakdash}{\nobreak\hbox{-}}
\setlength{\emergencystretch}{2em}
\multlinegap0pt
\usepackage{amssymb}
\usepackage{bm}
\usepackage{mathtools}
\usepackage{tikz}
\usepackage[shortlabels]{enumitem}
\newtheorem{lemma}{Lemma}
\newcommand{\ity}{\infty}
\title{Critical thresholds for the damped wave system with mixed product-type nonlinearities}
\author{Trung Loc Tang, Tuan Anh Dao, The Anh Cung}
\date{Source: arXiv:2609.21975v1 (CC BY 4.0)}
\begin{document}
\maketitle

\noindent\emph{Source and licence.} Critical thresholds for the damped wave system with mixed product-type nonlinearities. Version arXiv:2609.21975v1 (CC BY 4.0), distributed by arXiv under the Creative Commons Attribution 4.0 International licence, http://creativecommons.org/licenses/by/4.0/, as recorded in the arXiv licence metadata for this exact version and shown on the abstract page https://arxiv.org/abs/2609.21975v1. Full licence notice: \texttt{licenses/arxiv-2609.21975-CC-BY-4.0.txt}.\par\smallskip

\noindent\textit{Editorial scope.} Counted unit: the article's lemma regularity\_K (source0.tex lines 1875--1899). The article's complete original proof of it is delivered with it (source0.tex lines 1900--1985, one \texttt{proof} environment of 86 lines). No mathematics is written, repaired or re-derived: the statement and the proof are the article's own lines, in the article's own notation, and no retained line is substituted or re-flowed.  No further source-local context is needed: every cross-reference of every retained line resolves inside the two delivered spans.  Nothing was omitted from any retained span, so the package carries no omission record.  No retained line contains a citation, so the wrapper carries no bibliography and no bibliography entry.  No retained line contains a figure environment, an image-inclusion command or drawing code, so no image is delivered and the package declares no asset dependency.  Editorial formatting only, in addition: an AMS \texttt{amsart} preamble that restates the article's own declarations and macros used by the retained lines, namely the environments \texttt{lemma} on separate counters, together with the standard packages those lines need.  The retained environments print in this document as Lemma 1 in order of appearance, whereas the article prints the same items with the numbering of its own full document.  The complete native source of the article as distributed in the arXiv e-print for this exact version is bundled unchanged as \texttt{sources/211004/source0.tex}.
	\begin{lemma}\label{regularity_K}
		Let $T>0$ and $g\in L^1([0,T];L^2)$. We define
		\[
		v(t,\cdot):=\int_0^t\mathcal K(t-s,\cdot)\ast g(s,\cdot)\,ds,
		\qquad t\in[0,T].
		\]
		Then, $v\in \mathcal{C}([0,T];H^1)\cap \mathcal{C}^1([0,T];L^2)$ and $v$ is a weak solution to 
		\[
		\left\{
		\begin{aligned}
			&v_{tt}-\Delta v+ v_t= g,&&x\in \mathbb R^n,t\in (0,T),\\
			&v(0,x)=0, \quad v_t(0,x)=0,&&x\in \mathbb R^n.
		\end{aligned}
		\right.
		\]
		Furthermore, for any
		$\Phi\in \mathcal{C}^\infty_{\rm c}([0,T)\times\mathbb R^n)$ it holds
		\begin{equation}\label{duhamel_distribution_identity}
			\int_0^T\int_{\mathbb R^n}\Bigl(-\partial_tv\,\partial_t\Phi+\nabla v\cdot\nabla\Phi+\partial_tv\,\Phi- g\,\Phi\Bigr)\,dx\,dt=0
		\end{equation}
		and we have the following estimate:
		\begin{equation}\label{Energy_Est}
			\sup_{0\leq t\leq T}\big\{\|\partial_tv(t,\cdot)\|_{L^2}+\|\nabla v(t,\cdot)\|_{L^2}\big\}\lesssim \|g\|_{L^1([0,T];L^2)}.
		\end{equation}
	\end{lemma}

	\begin{proof}
		For $h\in L^2$, set $w(t,\cdot)=\mathcal K(t,\cdot)\ast h(\cdot)$ and $\partial_tw(t,\cdot)=\partial_t\mathcal K(t,\cdot)\ast h(\cdot)$. By the definition of $\mathcal K(t,\cdot)$ and the formula of Parseval-Plancherel, it holds
		\begin{equation}\label{uniformbound_H_1_and_t}
			\sup_{0\leq t\leq T}\big\{\|\partial_tw(t,\cdot)\|_{L^2}+\|\nabla w(t,\cdot)\|_{L^2}\big\}\lesssim \|h(\cdot)\|_{L^2}.
		\end{equation}
		Since $w(0,\cdot)=0$, we get $\|w(t,\cdot)\|_{L^2}\leq t\|h(\cdot)\|_{L^2}$. Thus, for any $T>0$ one has
		\begin{equation}\label{regularity_K_eq1}
			\sup_{0\leq t\leq T}\big\{\|w(t,\cdot)\|_{H^1}+\|\partial_tw(t,\cdot)\|_{L^2}
			\big\}\leq C(T)\|h(\cdot)\|_{L^2}.
		\end{equation}
		Moreover, the Fourier multipliers of $\mathcal{K}(t,\cdot)$ and $\partial_t K(t,\cdot)$ are continuous in $t$. So, we conclude that
		\begin{equation}\label{regularity_K_eq2}
			w\in \mathcal{C}([0,T];H^1) \quad \text{ and }\quad
			\partial_tw \in \mathcal{C}([0,T];L^2).
		\end{equation}
		We now write
		\[
		v(t,\cdot)=\int_0^t\mathcal{K}(t-s,\cdot)\ast g(s,\cdot)\,ds.
		\]
		In the view of \eqref{regularity_K_eq1}, this Bochner integral is well-defined and we may conclude the estimate \eqref{Energy_Est}. Let us next prove \eqref{duhamel_distribution_identity}. By \eqref{regularity_K_eq1}, we also have
		\[
		\|v(t,\cdot)\|_{H^1}
		\leq C(T)\int_0^t \|g(s,\cdot)\|_{L^2}\,ds.
		\]
		We next prove the continuity of $v$. Let $\tau>0$ be such that $t+\tau\leq T$. Then, it holds
		\[
		\begin{aligned}
			v(t+\tau,\cdot)-v(t,\cdot) &=\int_0^t
			\bigl(\mathcal{K}(t+\tau-s,\cdot)-\mathcal{K}(t-s,\cdot)\bigr)\ast g(s,\cdot)\,ds \\
			&\qquad +\int_t^{t+\tau}\mathcal{K}(t+\tau-s,\cdot)\ast g(s,\cdot)\,ds.
		\end{aligned}
		\]
		For a.e $s\in(0,t)$, the first integrand converges to zero in $H^1$ as $\tau\to0$, by
		\eqref{regularity_K_eq2}. Moreover, by
		\eqref{regularity_K_eq1}, it is bounded by $2C(T)\|g(s,\cdot)\|_{L^2}$, which belongs to $L^1(0,t)$. Hence, the first integral tends to zero by the dominated convergence theorem. The second integral satisfies
		\[
		\left\|\int_t^{t+\tau}\mathcal{K}(t+\tau-s,\cdot)\ast g(s,\cdot)\,ds\right\|_{H^1}\leq C(T)\int_t^{t+\tau}\| g(s,\cdot)\|_{L^2}\,ds,
		\]
		which tends to zero as $\tau \to 0$. The case $\tau<0$ is analogous. Therefore, $v\in \mathcal{C}([0,T];H^1)$. Taking the time-derivative of the Bochner convolution gives
		\begin{equation}\label{regularity_K_eq4}
			\partial_tv(t,\cdot)=\int_0^t\partial_t \mathcal{K}(t-s,\cdot)\ast g(s,\cdot)\,ds.
		\end{equation}
		Indeed, we have the following convergence for the first term 
		\[
		\int_0^t
		\frac{\bigl(\mathcal{K}(t+\tau-s,\cdot)-\mathcal{K}(t-s,\cdot)\bigr)\ast g(s,\cdot)}{\tau}\,ds\to \int_0^t\partial_t \mathcal{K}(t-s,\cdot)\ast g(s,\cdot)\,ds\text{ in }L^2
		\]
		as $\tau\to 0$ due to the dominated convergence theorem while the second term vanishes because
		\[
		\left\|\frac{1}{\tau}\int_t^{t+\tau}\mathcal{K}(t+\tau-s,\cdot)\ast g(s,\cdot)\,ds \right\|_{L^2}\leq \int_t^{t+\tau}\|g(s,\cdot)\|_{L^2}\,ds\to 0
		\]
		as $\tau\to 0$. Finally, by using \eqref{regularity_K_eq4}, for $\tau>0$ we get
		\[
		\begin{aligned}
			\partial_tv(t+\tau,\cdot)-\partial_tv(t,\cdot)&=\int_0^t
			\bigl(\partial_t \mathcal{K}(t+\tau-s,\cdot)-\partial_t\mathcal{K}(t-s,\cdot)\bigr)\ast g(s,\cdot)\,ds\\
			&\qquad +\int_t^{t+\tau}\partial_t \mathcal{K}(t+\tau-s,\cdot)\ast g(s,\cdot)\,ds.
		\end{aligned}
		\]
		The first integral tends to zero in $L^2$ by \eqref{regularity_K_eq1}, \eqref{regularity_K_eq2}, and the dominated convergence theorem, whereas
		\[
		\left\|\int_t^{t+\tau}\partial_t\mathcal{K}(t+h-s,\cdot)g(s,\cdot)\,ds\right\|_{L^2}\leq C(T)\int_t^{t+\tau}\| g(s,\cdot)\|_{L^2}\,ds\to 0.
		\]
		Thus, $\partial_tv\in \mathcal{C}([0,T];L^2)$. In particular, the formulas above give $v(0,\cdot)=0, \partial_tv(0,\cdot)=0$. It remains to verify the distributional identity. After choosing the function $g_\ell\in \mathcal{C}^\infty_{\rm c}((0,T)\times\mathbb R^n)$ such that $g_\ell\to g$ in $L^1([0,T];L^2)$ as $\ell\to\ity$, we define
		\[
		v_\ell(t,\cdot):=\int_0^t\mathcal{K}(t-s,\cdot)g_\ell(s,\cdot)\,ds.
		\]
		For any $\ell\in\mathbb N$, $v_\ell\in \mathcal{C}([0,T],H^2)\cap \mathcal{C}^1([0,T],H^1)\cap \mathcal{C}^2([0,T],L^2)$ is a strong solution to
		\[
		\left\{
		\begin{aligned}
			&\partial_t^2v_\ell-\Delta v_\ell+\partial_tv_\ell=g_\ell,&&x\in \mathbb{R}^n,t>0,\\
			&v_\ell(0,\cdot)=0,\quad \partial_tv_\ell(0,\cdot)=0,&&x\in \mathbb{R}^n.
		\end{aligned}
		\right.
		\]
		Hence, for any $\Phi\in \mathcal{C}^\infty_{\rm c}([0,T)\times\mathbb R^n)$, the integral-by-part fomurla implies that
		\begin{equation}\label{weak-identity}
			\int_0^T\int_{\mathbb R^n}\Bigl(-\partial_tv_\ell\,\partial_t\Phi+\nabla v_\ell\cdot\nabla\Phi+\partial_tv_\ell\,\Phi-g_\ell\Phi\Bigr)\,dx\,dt=0.
		\end{equation}
		By \eqref{regularity_K_eq1} and \eqref{regularity_K_eq4}, we arrive at
		\[
		\sup_{0\leq t\leq T}\Bigl(\|v_\ell(t,\cdot)-v(t,\cdot)\|_{H^1}+\|\partial_tv_\ell(t,\cdot)-\partial_tv(t,\cdot)\|_{L^2} \Bigr)\leq C(T)\|g_\ell- g\|_{L^1([0,T];L^2)} \to 0
		\]
		as $\ell\to\ity$. Passing to the limit in the preceding weak identity \eqref{weak-identity} yields \eqref{duhamel_distribution_identity}. Therefore, $v$ is a distributional solution with zero initial data. This completes the proof.
	\end{proof}

\end{document}
